% Options for packages loaded elsewhere
\PassOptionsToPackage{unicode}{hyperref}
\PassOptionsToPackage{hyphens}{url}
\PassOptionsToPackage{dvipsnames,svgnames,x11names}{xcolor}
\documentclass[
  11pt,
]{article}
\usepackage{xcolor}
\usepackage[margin=27mm]{geometry}
\usepackage{amsmath,amssymb}
\setcounter{secnumdepth}{-\maxdimen} % remove section numbering
\usepackage{iftex}
\ifPDFTeX
  \usepackage[T1]{fontenc}
  \usepackage[utf8]{inputenc}
  \usepackage{textcomp} % provide euro and other symbols
\else % if luatex or xetex
  \usepackage{unicode-math} % this also loads fontspec
  \defaultfontfeatures{Scale=MatchLowercase}
  \defaultfontfeatures[\rmfamily]{Ligatures=TeX,Scale=1}
\fi
\usepackage{lmodern}
\ifPDFTeX\else
  % xetex/luatex font selection
  \setmainfont[]{Cambria}
  \setmonofont[]{Consolas}
\fi
% Use upquote if available, for straight quotes in verbatim environments
\IfFileExists{upquote.sty}{\usepackage{upquote}}{}
\IfFileExists{microtype.sty}{% use microtype if available
  \usepackage[]{microtype}
  \UseMicrotypeSet[protrusion]{basicmath} % disable protrusion for tt fonts
}{}
\makeatletter
\@ifundefined{KOMAClassName}{% if non-KOMA class
  \IfFileExists{parskip.sty}{%
    \usepackage{parskip}
  }{% else
    \setlength{\parindent}{0pt}
    \setlength{\parskip}{6pt plus 2pt minus 1pt}}
}{% if KOMA class
  \KOMAoptions{parskip=half}}
\makeatother
\setlength{\emergencystretch}{3em} % prevent overfull lines
\providecommand{\tightlist}{%
  \setlength{\itemsep}{0pt}\setlength{\parskip}{0pt}}
\usepackage{bookmark}
\IfFileExists{xurl.sty}{\usepackage{xurl}}{} % add URL line breaks if available
\urlstyle{same}
\hypersetup{
  pdftitle={WP-11 · Formalism: the interleaved bilateral construction and its two properties},
  pdfauthor={The Privacy-is-Value Research Programme},
  colorlinks=true,
  linkcolor={black},
  filecolor={Maroon},
  citecolor={Blue},
  urlcolor={blue},
  pdfcreator={LaTeX via pandoc}}

\title{WP-11 · Formalism: the interleaved bilateral construction and its
two properties}
\usepackage{etoolbox}
\makeatletter
\providecommand{\subtitle}[1]{% add subtitle to \maketitle
  \apptocmd{\@title}{\par {\large #1 \par}}{}{}
}
\makeatother
\subtitle{WP-11 (rpp-adversarial-paper) · tier A (formal sections for
the tier-A paper) · 2026-07-18 · status DRAFT v1 --- definitions + one
conditional reduction; every st}
\author{The Privacy-is-Value Research Programme}
\date{2026-07-18}

\begin{document}
\maketitle

This section formalises the two properties the RPP construction claims
under the First-Person ruling (Option C): a \emph{behavioral} property
of the injection mechanism (forcing), and an
\emph{information-theoretic} property of the two-party interleaved pair
(withheld-fraction security). The central discipline is that these are
separate: forcing gets the proverb authored; the entropy split is what
makes the result secure. Conflating them --- presenting the behavioral
property as a security guarantee --- is the failure mode this section
exists to prevent (the modal error of E6-C06: a mechanism that compels
is not thereby a mechanism that cannot be broken).

\subsection{1. The construction}\label{the-construction}

Two parties A and B share an experience E and each author a
natural-language compression P\_A, P\_B (their proverbs). An
\textbf{interleave} map I combines them into a single object I(P\_A,
P\_B); a public \textbf{anchor} discloses a σ-fraction of that object,
A\_σ := disc\_σ(I(P\_A, P\_B)) for a disclosure ratio σ ∈ {[}0,1{]},
while the complementary \textbf{withheld fraction} W := I(P\_A, P\_B)
~A\_σ (the (1−σ)-fraction) is held jointly by the two parties. The
commitment is κ := H(I(P\_A, P\_B)) for a collision-resistant hash H.
Verification/recovery requires reconstructing I(P\_A, P\_B) to within a
fuzzy tolerance sufficient to recompute κ (equivalently, to open the
commitment under an error-correcting opening; the hard-crypto layer ---
H, and the signatures/κ-addressing of the deployed system --- is assumed
from the implementation record E11 and is not re-derived here).

Deployment status (GR-8, E3-C32): I and the σ-budget are design-layer
objects (the deployed system uses the asymmetric-default inscription and
a stratum-derived visibility, not σ ∈ {[}0,1{]}); this section
formalises the specified construction and states this boundary.

\subsection{2. Forcing (the behavioral property of the injected
directive) --- Definition, not a
guarantee}\label{forcing-the-behavioral-property-of-the-injected-directive-definition-not-a-guarantee}

Let a reader R (a human or a model) process content C, and let Φ(C)
denote the event that R emits a proverb \emph{bound to C} before
producing any downstream answer. Let D be the injected RPP directive.
Let 𝒞 be a content distribution and 𝓡 a reader population.

\textbf{Definition 1 (forcing).} The directive D is \textbf{(f, b,
s)-forcing} over (𝒞, 𝓡) if 1. \emph{(compulsion)} E\_\{C∼𝒞, R∼𝓡\}{[}
Pr(Φ(C+D)) − Pr(Φ(C)) {]} ≥ f (the directive raises the probability of
the bound act by at least f over the no-directive baseline); 2.
\emph{(content-binding)} Pr( the emitted proverb is specific to C
\textbar{} Φ(C+D) ) ≥ b, where ``specific to C'' is failure of a fixed
generic-completion test (a proverb that a reader who never saw C could
also have produced does not count); 3. \emph{(strip-survival)} under
each of a declared family of pre-processing attacks 𝒜 (directive
deletion, paraphrase, relocation), the compulsion gap remains ≥ s.

\textbf{Remark 1 (what Definition 1 is and is not).} Forcing is a
statistical property of a reader population under a directive. It is NOT
unforgeability and NOT a security guarantee: a reader that is forced to
author a bound proverb has done exactly what the mechanism intends, and
a capable model performing the act is a success of the mechanism, not a
break of it. The security of the construction does not rest on
Definition 1 at all (it rests on §3). The only role of forcing is to
ensure the proverb is authored under the directive rather than skipped.

\textbf{Remark 2 (forcing is the dual of injection-resistance, hence
model-relative).} Prompt-injection \emph{attacks} seek to make a model
follow an injected instruction; prompt-injection \emph{defenses} seek
resistance (Greshake et al.~arXiv:2302.12173; OWASP LLM01). RPP's
forcing is cooperative injection: it wants the model to follow.
Therefore f is inversely related to the reader's injection-resistance,
and the honest claim is bounded --- forcing holds over compliant reader
populations and weakens on hardened ones. The strongest positive form of
the claim is that D is (f, b, s)-forcing even over an
injection-resistant population (the benchmark's arm A3, L164); the
weakest admissible form scopes 𝓡 to cooperative readers and says so.
Definition 1 carries 𝓡 explicitly so the scope is never hidden.

\textbf{Distinction from neighbours, in the definition's terms.}
Zero-shot chain-of-thought (Kojima et al.~arXiv:2205.11916) achieves
compulsion of a reasoning act (f high) with content-binding b ≈ 0 (the
reasoning is not a C-specific authored artifact bound to a commitment)
and is self-directed (the operator steers their own model, not a
directive embedded in third-party content). ``Context bombs'' (Tracebit
2026) raise a \emph{derailment} event, not Φ, over a hostile reader. RPP
is the (high f, high b, embedded-and-cooperative) corner; b is the
coordinate that separates it from CoT, and the cooperative-embedded
setting separates it from context bombs.

\subsection{3. Withheld-fraction security (the information-theoretic
property of the
pair)}\label{withheld-fraction-security-the-information-theoretic-property-of-the-pair}

Fix the malicious-counterparty threat model (the honest adversary,
L163): the adversary is party B, who was present at E and now holds
their own proverb P\_B, the shared-experience context E, and the public
anchor A\_σ. Write the adversary's side information S := (P\_B, E,
A\_σ). The adversary wins by producing W′ within fuzzy tolerance t of
the withheld fraction W (enough to recompute κ). We use fuzzy
conditional min-entropy in the sense of the fuzzy-extractor line (Dodis,
Reyzin \& Smith, EUROCRYPT 2004; Juels \& Wattenberg, CCS 1999):
H̃\emph{\{t,∞\}(X \textbar{} Z) := −log E}\{z\}{[} max\_\{x′\} Pr( d(X,
x′) ≤ t \textbar{} Z = z ) {]}, the log of the best t-ball guessing
probability given Z.

\textbf{Definition 2 (interleave splitting).} The interleave I is
\textbf{(σ, δ)-splitting} if the σ-disclosure leaks at most δ bits about
the withheld fraction: I(A\_σ ; W) ≤ δ. (δ = 0 is an ideal interleave
whose public fraction is independent of the withheld fraction; a poor
interleave, e.g.~a prefix cut of a redundant encoding, has large δ and
is insecure regardless of proverb entropy. δ is a property of the map I,
measurable directly.)

\textbf{Definition 3 (informed-counterparty security).} The construction
is \textbf{(t, ε)-secure} against the malicious counterparty if the
guessing probability of the withheld fraction within tolerance t given
the adversary's side information is at most ε: \textgreater{} Pr{[} A(S)
outputs W′ with d(W′, W) ≤ t {]} ≤ ε, equivalently H̃\_\{t,∞\}(W
\textbar{} S) ≥ log(1/ε).

\textbf{Proposition 4 (security reduction --- Proven-conditional).} If I
is (σ, δ)-splitting (Definition 2) and party A's proverb retains fuzzy
conditional min-entropy H̃\emph{\{t,∞\}(P\_A \textbar{} P\_B, E) ≥ m
against a co-participant who holds P\_B and E, then the construction is
(t, ε)-secure with \textgreater{} ε ≤ 2\^{}\{−(m − δ)\}. \emph{Proof.}
The withheld fraction W is a deterministic function of (P\_A, P\_B) via
I. The adversary's guess of W within tolerance t is at least as hard as
guessing the part of I(P\_A, P\_B) not determined by (P\_B, E) and not
disclosed by A\_σ. Conditioning reduces min-entropy by at most the
information the conditioning carries: H̃}\{t,∞\}(W \textbar{} P\_B, E,
A\_σ) ≥ H̃\emph{\{t,∞\}(W \textbar{} P\_B, E) − I(A\_σ ; W \textbar{}
P\_B, E) ≥ (m) − δ, where the first inequality is the standard chain
bound for conditional min-entropy under side information of at most
I(A\_σ;W\textbar P\_B,E) ≤ δ bits (Definition 2, taken conditionally),
and the second uses that the residual entropy of W given (P\_B, E) is
lower-bounded by that of P\_A given (P\_B, E) up to the fuzzy tolerance,
since W includes the (1−σ)-fraction of A's contribution. The guessing
bound H̃}\{t,∞\}(W \textbar{} S) ≥ m − δ gives ε ≤ 2\^{}\{−(m−δ)\}. ∎
\textbf{Named gaps (A3 discipline --- do not hand-wave).} Two steps are
stated as requirements the construction must meet, not as facts: (i) the
claim that W's residual entropy given (P\_B, E) is lower-bounded by
P\_A's requires that the interleave place at least the withheld share of
A's independent contribution into W --- an
\textbf{interleave-faithfulness requirement} on I, listed as an
engineering requirement, not assumed; (ii) the fuzzy-tolerance
accounting (that error-correcting to within t does not itself leak more
than the secure-sketch entropy loss) is the standard fuzzy-extractor
caveat and is inherited, with its known entropy-loss term folded into m
(i.e.~m is the \emph{post-sketch} residual entropy). Both are
inherited-or-required, and both are measurable.

\textbf{Corollary 5 (the estimator link).} The benchmark quantity r(σ)
(L164, H2) is an empirical estimator of ε(σ). The pre-registered
falsification r(0.618) ≥ 0.5 is exactly the formal statement m − δ ≤ 1
bit at σ = 0.618: the withheld majority-public sweet-spot retains at
most one bit of residual entropy against an informed counterparty,
i.e.~the construction is insecure there. Thus the experiment measures
the two quantities Proposition 4 reduces security to (δ from I, m from
proverb entropy given a co-participant), and refutes or supports the
reduction's premises directly.

\textbf{Proposition 6 (monotonicity in σ --- Conjecture, proof sketch +
missing step).} If I is \textbf{graceful} --- each fragment moved from
the disclosed set A\_σ to the withheld set W as σ decreases adds
non-negative conditional min-entropy to W given S --- then ε(σ) is
non-increasing as σ decreases, so r(σ) is non-increasing in σ⁻¹ (more
withheld ⇒ harder). \emph{Proof sketch.} Decreasing σ enlarges W by
fragments previously in A\_σ; under gracefulness each such fragment
weakly increases H̃\_\{t,∞\}(W \textbar{} S); the guessing bound of
Proposition 4 then gives non-increasing ε. \emph{Missing step (named).}
Gracefulness is not automatic: an interleave can move a fragment into W
that is fully determined by fragments remaining in A\_σ (or by P\_B),
adding zero entropy, or worse, changing the correlation structure so δ
rises. Whether a concrete I is graceful is a property of I to be
\textbf{constructed and proven}, not assumed; absent that, monotonicity
is a conjecture the benchmark's isotonic test (L164) can refute
empirically. This is labelled Conjecture, not Proposition-with-∎.

\subsection{4. The layering (the honesty discipline this section
enforces)}\label{the-layering-the-honesty-discipline-this-section-enforces}

The two properties occupy different registers and must be stated as such
throughout the paper: - \textbf{Forcing (Definition 1)} is behavioral,
model-relative, and provides no security. Its role is to author the
proverb under the directive. It is measured by H1. - \textbf{Security
(Definitions 2--3, Proposition 4)} is information-theoretic and holds
regardless of whether any reader ``understood'': it rests on the
residual entropy of A's proverb given a co-participant's full knowledge,
and on the interleave's splitting. It is conditional on two measurable
quantities (m, δ) and on two named engineering requirements
(interleave-faithfulness, gracefulness). It is measured by H2.

\textbf{Consequence for the brain-wallet discharge (making L163
precise).} The brain-wallet critique is the statement that a single
human-memorable secret has low min-entropy against a public target.
Proposition 4 shows the security here is not the min-entropy of one
proverb but H̃\_\{t,∞\}(P\_A \textbar{} P\_B, E) --- the residual entropy
of A's proverb \emph{given a co-participant's complete knowledge of the
shared experience and their own half}. The critique is therefore not
refuted by assertion but relocated to a measurable quantity: if a
co-participant can guess A's proverb (m small), the construction is
insecure, and H2 measures exactly this. The honest threat is the
informed co-participant, and the security claim is precisely ``m − δ is
large at the operating σ,'' neither more nor less.

\textbf{Coherence with the programme (cited, not conflated).} Definition
3 and Proposition 4 are the same information-theoretic family as the
companion's informed-adversary corollary (WP-07 Corollary 5.4b: an error
floor governed by H(X \textbar{} B\_t) against an adversary holding
accumulating side information). Both bound reconstruction by a residual
conditional entropy against an informed adversary. They are not the same
theorem --- WP-07's B\_t is accumulating linkage over calendar time
about a fixed archive; here S is a fixed co-participant's knowledge of a
shared experience --- and the paper cites the kinship without asserting
identity (GR-1/GR-4 discipline: same family, distinct objects).

\subsection{5. Firewall --- what is proved, defined, assumed, empirical,
conjectural}\label{firewall-what-is-proved-defined-assumed-empirical-conjectural}

\begin{itemize}
\tightlist
\item
  \textbf{Proved-conditional:} Proposition 4 (the reduction ε ≤
  2\^{}\{−(m−δ)\}), by standard conditional-min-entropy bounds,
  conditional on Definitions 2 and the two named requirements.
\item
  \textbf{Definitions:} Definition 1 (forcing), Definitions 2--3
  (splitting, security). Forcing is a definition whose
  \emph{satisfaction} is empirical.
\item
  \textbf{Empirical (measured, not assumed):} m (residual proverb
  entropy given a co-participant --- H2), δ (interleave leakage ---
  direct measurement), f/b/s (forcing --- H1), the σ-frontier crossover
  (H2×H3).
\item
  \textbf{Engineering requirements (must be met by the construction,
  listed not assumed):} interleave-faithfulness (Prop 4 gap i),
  gracefulness (Prop 6), the fuzzy-sketch entropy-loss accounting (Prop
  4 gap ii).
\item
  \textbf{Conjecture:} Proposition 6 (monotonicity in σ), with the
  missing step named (gracefulness of a concrete I).
\item
  \textbf{Cited substrate, not claimed:} the
  fuzzy-extractor/secure-sketch primitives (Dodis-Reyzin-Smith 2004;
  Juels-Wattenberg 1999) --- RPP builds on them for the
  reproduce-from-noisy-input layer and does not claim that primitive
  (WP-11a).
\item
  \textbf{Conjectural parameters:} the φ candidate values for σ
  (register C54) --- candidate operating points, never asserted optima.
\item
  \textbf{Assumed from elsewhere:} the hard-crypto layer (H, κ,
  signatures) from E11; not re-derived.
\end{itemize}

Nothing in this section asserts the construction unconditionally secure.
The strongest statement is: \emph{the construction is (t,
2\^{}\{−(m−δ)\})-secure against an informed co-participant, where m and
δ are quantities the pre-registered experiment measures, under two named
interleave requirements; whether m − δ is large enough at an operating σ
is exactly what H2 decides.}

\subsection{6. The concrete interleave I* (First-Person ruling,
2026-07-18) --- the construction that
binds}\label{the-concrete-interleave-i-first-person-ruling-2026-07-18-the-construction-that-binds}

The abstract interleave I is now fixed to a concrete construction, so
Propositions 4 and 6 have an object to bind to. The choice was between a
graded interleave (σ a smooth disclosure knob) and an all-or-nothing
transform (maximal security but σ binary); the First Person ruled the
graded path, which matches the ``interleaved visibility ratio in this
knowledge sharing'' intent (L163) and keeps σ a real tunable knob.

\textbf{Definition 7 (the concrete interleave I*,
extract-then-block-interleave).} Fix a fuzzy extractor (secure sketch SS
+ strong extractor Ext) with public helper data and error tolerance t
(Dodis-Reyzin-Smith, EUROCRYPT 2004). Then: 1. \textbf{Extract.} s\_A :=
Ext(SS(P\_A)), s\_B := Ext(SS(P\_B)); each is (near-)uniform of length
equal to its post-sketch residual min-entropy; helpers are public. This
step does two jobs: it tolerates rewording (the fuzzy part) and it
whitens natural-language redundancy so blocks become (near-)independent
and uniformly informative. 2. \textbf{Block-interleave.} Partition s\_A,
s\_B into n blocks each and interleave on a fixed public symmetric
schedule π (e.g.~s\_A{[}1{]} s\_B{[}1{]} s\_A{[}2{]} s\_B{[}2{]}
\ldots), giving Z := π(s\_A, s\_B). 3. \textbf{Disclose.} The anchor
A\_σ is a σ-fraction of Z's blocks under a public schedule that draws
the same fraction from each party's blocks; W is the complement; κ :=
H(Z).

\textbf{What I* discharges (the two named gaps of §3, reduced to the
extractor --- not assumed).} - \emph{Splitting (Definition 2), δ bounded
not merely measured.} Ext output is within statistical distance δ\_ext
of uniform, and s\_A ⟂ s\_B by independent authoring (the N2 property).
A co-participant B re-derives the s\_B blocks from P\_B and the public
helper, so the s\_B blocks in A\_σ leak nothing new to B; the s\_A
blocks are (near-)uniform, so disclosing a σ-fraction of them leaks at
most δ\_ext about the withheld (1−σ). Hence I(A\_σ ; W \textbar{} P\_B,
E) ≤ δ\_ext, and δ is a bounded extractor parameter, not an empirical
unknown. - \emph{Interleave-faithfulness (Proposition 4, gap i),
DISCHARGED.} The symmetric schedule places (1−σ) of s\_A's blocks into W
by construction; since s\_A carries A's post-sketch residual entropy
against B, W's residual entropy given the adversary's side information
is at least (1−σ)·m up to δ\_ext, where m := H̃\emph{\{t,∞\}(P\_A
\textbar{} P\_B, E) post-sketch. The requirement is met by the schedule,
not assumed. - \emph{Gracefulness (Proposition 6), PROVEN for I*.} Each
block moved from A\_σ to W as σ decreases is a (near-)uniform block
(near-)independent of the disclosed set, so it adds (near-)its full
entropy to H̃}\{t,∞\}(W \textbar{} S). Monotonicity therefore holds for
I* --- Proposition 6, a Conjecture for arbitrary I, is a Proposition for
I*: ε(σ) is non-increasing as σ decreases, strictly under the uniformity
of Ext.

\textbf{Theorem 8 (security of I*, Proven-conditional on the extractor
and on m).} Under Definition 7, the construction is (t, ε)-secure
against the malicious informed counterparty with \textgreater{} ε ≤
2\^{}\{−((1−σ)·m − δ\_ext)\}, m := H̃\emph{\{t,∞\}(P\_A \textbar{} P\_B,
E) post-sketch, δ\_ext the extractor's uniformity loss. \emph{Proof.}
Proposition 4 with the I*-discharged inputs: faithfulness gives
H̃}\{t,∞\}(W \textbar{} S) ≥ (1−σ)·m − δ\_ext; the guessing bound gives ε
≤ 2\^{}\{−((1−σ)·m − δ\_ext)\}. Monotonicity is Proposition 6 for I*. ∎
\textbf{The single remaining empirical unknown.} Choosing I* collapses
the two engineering requirements and the δ-measurement into the
extractor's standard guarantee (δ\_ext bounded, faithfulness by
schedule, gracefulness proven). Exactly one quantity remains empirical:
\textbf{m, the residual entropy of a proverb given a co-participant who
was present at the shared experience}. The entire security question is
now ``how guessable is A's proverb to someone who was there,'' which is
the honest threat and is exactly what benchmark H2 measures. A low-m
proverb is insecure under any interleave; I* guarantees you retain
(1−σ)·m of whatever m is, no less.

\textbf{The two rejected alternatives, recorded.} - \emph{Naive text
interleave (no extractor):} interweave the raw proverb characters/words.
Human-readable, but natural-language redundancy makes disclosed
fragments predict withheld ones (δ large), so security is weak and not
cleanly analysable. Rejected as the security layer; it may serve only as
the readable surface (below). - \emph{All-or-nothing transform (AONT):}
withhold any fragment ⇒ the whole is unrecoverable (δ = 0, maximal
security), but σ collapses to binary and the graded knowledge-sharing
knob is lost. Rejected for the graded-σ intent; it remains the option if
maximal security is wanted and graded disclosure is not.

\textbf{Layering (keeps M2 honest).} I* whitens proverbs to bits for the
SECURITY layer (Z); the human-readable proverb is a SEPARATE surface
layer for knowledge-sharing and relational recall (M2/M3). σ applies to
both in lockstep: the readable anchor discloses σ of the joint
human-readable understanding, and the withheld (1−σ) of Z is what
protects. The two are coupled --- sharing more readable knowledge
(higher σ) genuinely lowers security by lowering (1−σ)·m --- which is
the σ-frontier the benchmark measures, not a cosmetic knob. The
secure-sketch helper data's own leakage is folded into m being the
post-sketch residual (the standard fuzzy-extractor accounting, inherited
from the cited primitive).

\textbf{Design rationale --- why the graded path (the
understanding-as-key thesis, restated register-neutrally).} The choice
of I* over the all-or-nothing transform is not merely a usability
preference: in the understanding-as-key research the disclosure-security
trade IS the intended generative mechanism, framed there as
\emph{identity as a dance, not a stance} (First-Person framing).
Register-neutrally: a stance is a static held secret verified by
equality; a dance is an ongoing relational disclosure process in which
parties continuously trade what they reveal against what they hold, and
the trajectory of that trade over time --- not any fixed value --- is
the identity. The graded σ of I* is the formal shadow of that thesis: it
makes disclosure a continuous, re-negotiable position (a dance along σ),
where AONT would force a static all-or-nothing secret (a stance). The
residual-entropy kernel is already present in the understanding-as-key
source: ``attackers cannot enumerate without understanding the
relationship itself'' (understanding-as-key :872) is exactly m =
H̃\_\{t,∞\}(P\_A \textbar{} P\_B, E) --- security as the residual
uncertainty of the relationship given the adversary's own knowledge of
it. GR-4/GR-8 provenance (located 2026-07-18): the phrasing ``identity
is the dance, not the stance'' is in the corpus at
poems/tide-orbit-selene.md:26 (with its development at :101-108, ``the
dance is the identity''), which is NARRATIVE-CANON (the seventh-capital
/ C55 home, per SOURCES seventh-capital-canon). Per GR-4 narrative-canon
is restated register-neutrally and NOT cited directly in a tier-A
artifact; the citable formal anchor is the understanding-as-key/Zypher
kernel at :872 (m = residual relationship-entropy). The Zypher lineage
paper carries the kernel but not the phrase; the phrase's home is the
poem. (The First Person recalls a Nov-2025 understanding-as-key document
also carrying it; if that is a source distinct from the poem it is a
minor SOURCES-completeness item for A1, but the concept is now grounded
and the formal statements stand independently of the framing
regardless.)

\subsection{7. Handoff}\label{handoff}

\begin{itemize}
\tightlist
\item
  \textbf{A6:} the definitions here fix what the metrics estimate ---
  r(σ) ↔ ε(σ) (Corollary 5), the forcing triple ↔ (f, b, s) (Definition
  1). No benchmark change required; the content-binding detector
  operationalises ``specific to C'' (Definition 1.2), and the
  interleave-leakage measurement operationalises δ (Definition 2). Add
  one measurement: δ(σ) directly (I(A\_σ; W)), so Proposition 4's two
  inputs are both measured, not one inferred.
\item
  \textbf{A5-pc (post-results):} the two places a PC breaks this are the
  named gaps in Proposition 4 (interleave-faithfulness) and the
  gracefulness conjecture (Proposition 6). Both are labelled and both
  are empirically refutable, which is the defensible posture; the review
  should check that the paper never states Proposition 6 as proved.
\item
  \textbf{A4:} cite Dodis-Reyzin-Smith 2004 and Juels-Wattenberg 1999 at
  primary records (both already in the WP-11a spine); confirm the fuzzy
  conditional min-entropy notation matches the cited source before
  tier-A use.
\item
  \textbf{First Person:} the interleave I is currently abstract; a
  concrete I must be chosen and shown faithful+graceful for Propositions
  4 and 6 to bind. This is a construction decision (which interleave),
  separate from the earlier architecture ruling. Until an I is fixed,
  the security is a template, not a theorem.
\end{itemize}

\end{document}
